\documentclass[UTF8,a4paper,11pt,fontset=fandol]{ctexart}
\usepackage[margin=2.3cm]{geometry}
\usepackage{amsmath,amssymb}
\usepackage[hidelinks]{hyperref}
\setlength{\parindent}{0pt}
\setlength{\parskip}{5pt}
\allowdisplaybreaks
\newcommand{\E}{\mathbb E}
\newcommand{\Pp}{\mathbb P}
\newcommand{\norm}[1]{\left\lVert#1\right\rVert}
\title{HW3 中文详解}
\author{00103335：深度学习与强化学习}
\date{2026 年秋季}
\begin{document}
\pagestyle{plain}
\maketitle
本详解按原题顺序说明计算和论证依据。

同目录文件：原题 \texttt{hw3.pdf}，英文解答 \texttt{solution.pdf}。

全篇不额外假设噪声独立、同分布或对称。

\section{随机误差递推：四阶矩、概率阶与几乎处处收敛}
题设给出
\[
 \delta_{t+1}=(1-\lambda\eta_t)\delta_t+\eta_t\xi_t,
 \qquad\eta_t=\frac{a}{(t+t_0)^\alpha},\qquad\frac12<\alpha\leq1.
\]
其中 $\lambda,a>0$，$0<\lambda\eta_t\leq1$；若 $\alpha=1$，还要求
$2\lambda a>1$。$\mathcal F_t$ 表示时刻 $t$ 已知的信息，
$\delta_t$ 对这些信息可测，所以取 $\mathcal F_t$ 条件期望时，
$\delta_t$ 可以像已知量一样移到期望外。噪声满足
\[
 \E[\xi_t\mid\mathcal F_t]=0,
 \qquad\E[|\xi_t|^4\mid\mathcal F_t]\leq\sigma^4.
\]
记 $q_t=1-\lambda\eta_t\in[0,1]$、$v_t=\E|\delta_t|^4$。

\subsection*{预备：低阶噪声矩与可积性}
条件 H\"older 不等式说明，对于 $r=2,3$，
\[
 \E[|\xi_t|^r\mid\mathcal F_t]
 \leq\big(\E[|\xi_t|^4\mid\mathcal F_t]\big)^{r/4}
 \leq\sigma^r.
\]
特别地，条件二阶矩不超过 $\sigma^2$，条件三阶矩的绝对值不超过
$\E[|\xi_t|^3\mid\mathcal F_t]\leq\sigma^3$。
零条件均值仅能消去一次噪声项，不能据此消去三次噪声项。

还要先确认每个 $v_t$ 有限。由凸性，
\[
 |u+v|^4\leq(|u|+|v|)^4
 =16\left(\frac{|u|+|v|}{2}\right)^4
 \leq8(|u|^4+|v|^4).
\]
将 $u=q_t\delta_t$、$v=\eta_t\xi_t$ 代入，再取期望，得到
\[
 v_{t+1}\leq8q_t^4v_t+8\eta_t^4\E|\xi_t|^4
 \leq8v_t+8\eta_t^4\sigma^4.
\]
由于 $v_1<\infty$，归纳可知任意有限时刻的四阶矩都存在。
这也保证后续展开的各混合项可积。

\subsection*{(a) 推导四阶矩的收缩递推}
目标是：任取 $c\in(0,4\lambda)$，最终都能得到
$v_{t+1}\leq(1-c\eta_t)v_t+C_0\eta_t^3$。

\textbf{第一步：完整展开四次方。}
因为 $\delta_t$ 是实数，$|\delta_{t+1}|^4=\delta_{t+1}^4$。二项式公式给出
\begin{align*}
 (q_t\delta_t+\eta_t\xi_t)^4
 &=q_t^4\delta_t^4+4q_t^3\delta_t^3\eta_t\xi_t
   +6q_t^2\delta_t^2\eta_t^2\xi_t^2\\
 &\quad+4q_t\delta_t\eta_t^3\xi_t^3+\eta_t^4\xi_t^4.
\end{align*}
给定 $\mathcal F_t$ 后，第二项的条件期望为零。其余项用预备矩界，
并使用 $q_t,q_t^2\leq1$，得到
\begin{align*}
 \E[|\delta_{t+1}|^4\mid\mathcal F_t]
 &\leq q_t^4\delta_t^4+6\sigma^2\eta_t^2\delta_t^2
   +4\sigma^3\eta_t^3|\delta_t|+\sigma^4\eta_t^4.
\end{align*}
三阶交叉项可能是正的，所以这里取其绝对值上界，而不是令其为零。

\textbf{第二步：把一次误差项合并到二次误差项。}
由 $(u-v)^2\geq0$ 可得 $2uv\leq u^2+v^2$。
取 $u=\sigma\eta_t|\delta_t|$、$v=\sigma^2\eta_t^2$，再乘以 $2$，便有
\[
 4\sigma^3\eta_t^3|\delta_t|
 \leq2\sigma^2\eta_t^2\delta_t^2+2\sigma^4\eta_t^4.
\]
与原有各项合并后，
\[
 \E[|\delta_{t+1}|^4\mid\mathcal F_t]
 \leq q_t^4\delta_t^4+8\sigma^2\eta_t^2\delta_t^2+3\sigma^4\eta_t^4.
\]

\textbf{第三步：用少量四阶项吸收二阶项。}
对任意 $\varepsilon>0$，平方非负给出
\begin{align*}
 0&\leq\left(\sqrt{\varepsilon\eta_t}\,\delta_t^2
      -\frac{4\sigma^2\eta_t^{3/2}}{\sqrt\varepsilon}\right)^2\\
  &=\varepsilon\eta_t\delta_t^4-8\sigma^2\eta_t^2\delta_t^2
    +\frac{16\sigma^4}{\varepsilon}\eta_t^3.
\end{align*}
移项即得
\[
 8\sigma^2\eta_t^2\delta_t^2
 \leq\varepsilon\eta_t\delta_t^4
 +\frac{16\sigma^4}{\varepsilon}\eta_t^3.
\]
这样右侧只剩四阶误差和确定性的步长幂次：
\[
 \E[|\delta_{t+1}|^4\mid\mathcal F_t]
 \leq(q_t^4+\varepsilon\eta_t)\delta_t^4
 +\frac{16\sigma^4}{\varepsilon}\eta_t^3+3\sigma^4\eta_t^4.
\]

\textbf{第四步：明确选取常数。}
固定题目指定的 $c<4\lambda$，取 $\varepsilon=(4\lambda-c)/2>0$。
对 $x\in[0,1]$，
\[
 (1-x)^4=1-4x+6x^2-4x^3+x^4\leq1-4x+6x^2,
\]
因为最后两项为 $x^3(x-4)\leq0$。
取 $x=\lambda\eta_t$，则
\[
 q_t^4+\varepsilon\eta_t
 \leq1-4\lambda\eta_t+6\lambda^2\eta_t^2+\varepsilon\eta_t.
\]
$\eta_t\to0$，所以可以选定 $t_1$，使所有 $t\geq t_1$ 满足
$\eta_t\leq1$ 及 $6\lambda^2\eta_t\leq\varepsilon$。此时
\[
 q_t^4+\varepsilon\eta_t
 \leq1-(4\lambda-2\varepsilon)\eta_t=1-c\eta_t,
 \qquad\eta_t^4\leq\eta_t^3.
\]
最后对条件不等式取普通期望，利用全期望公式，得到
\[
 \boxed{v_{t+1}\leq(1-c\eta_t)v_t+C_0\eta_t^3},
 \qquad C_0=\left(\frac{16}{\varepsilon}+3\right)\sigma^4<\infty.
\]
$C_0$ 和 $t_1$ 可以依赖固定参数，但都不随 $t$ 改变。
当 $\sigma=0$ 时同样成立，因为证明没有除以 $\sigma$。

\subsection*{(b) 证明 $v_t\leq C\eta_t^2$}
这里不能简单把 (a) 的稳态数量级当成证明，因为目标上界
$C\eta_t^2$ 本身随时间下降。需要比较递推的收缩速度和这个上界的下降速度。

\textbf{第一步：计算目标上界的相邻差。}
定义
\[
 D_t=\frac{\eta_t^2-\eta_{t+1}^2}{\eta_t^3},
 \quad\text{于是}\quad\eta_{t+1}^2=\eta_t^2-D_t\eta_t^3.
\]
令 $s=t+t_0$。代入 $\eta_t=as^{-\alpha}$ 得
\begin{align*}
 D_t&=\frac{a^2s^{-2\alpha}-a^2(s+1)^{-2\alpha}}{a^3s^{-3\alpha}}\\
 &=\frac{s^\alpha}{a}\left[1-\left(1+\frac1s\right)^{-2\alpha}\right].
\end{align*}
对函数 $h(u)=(1+u)^{-2\alpha}$ 在 $[0,1/s]$ 使用中值定理，
存在 $\zeta_s\in(0,1/s)$，使
\[
 1-(1+1/s)^{-2\alpha}
 =\frac{2\alpha}{s}(1+\zeta_s)^{-2\alpha-1}.
\]
因此
\[
 D_t=\frac{2\alpha}{a}s^{\alpha-1}(1+\zeta_s)^{-2\alpha-1}
 \longrightarrow d_*=
 \begin{cases}0,&\alpha<1,\\2/a,&\alpha=1.\end{cases}
\]

\textbf{第二步：为什么 $\alpha=1$ 时需要额外条件。}
我们需要选择两个固定常数，使 $d_*<d<c<4\lambda$。
当 $\alpha<1$ 时 $d_*=0$，总能选到。
当 $\alpha=1$ 时，能够选到的充要条件是
\[
 \frac2a<4\lambda\quad\Longleftrightarrow\quad2\lambda a>1,
\]
正是题设额外给出的条件。因此取上述 $d,c$，并使用 (a) 对这个 $c$ 的结论。
再选充分大的整数 $N$，使对所有 $t\geq N$，
\[
 v_{t+1}\leq(1-c\eta_t)v_t+C_0\eta_t^3,
 \qquad D_t\leq d,\qquad c\eta_t\leq1.
\]
最后一条保证 $1-c\eta_t\geq0$，后面代入上界时不会改变不等号方向。

\textbf{第三步：选定覆盖初始时刻的常数并归纳。}
取
\[
 C\geq\max\left\{\frac{C_0}{c-d},\
                 \max_{1\leq t\leq N}\frac{v_t}{\eta_t^2}\right\}.
\]
有限个 $v_t$ 都有限，所有 $\eta_t>0$，而 $c-d>0$，所以 $C<\infty$。
此选择首先保证 $1\leq t\leq N$ 时结论成立。
假设某个 $t\geq N$ 时 $v_t\leq C\eta_t^2$，则
\begin{align*}
 v_{t+1}
 &\leq(1-c\eta_t)C\eta_t^2+C_0\eta_t^3\\
 &=C\eta_t^2-(Cc-C_0)\eta_t^3\\
 &\leq C\eta_t^2-Cd\eta_t^3\\
 &\leq C\eta_t^2-CD_t\eta_t^3=C\eta_{t+1}^2.
\end{align*}
第三行因为 $C(c-d)\geq C_0$，第四行因为 $D_t\leq d$。
归纳完成，得到对所有 $t$ 都成立的
\[
 \boxed{\E|\delta_t|^4\leq C\eta_t^2.}
\]

\subsection*{(c) 从四阶矩推出概率阶}
$\delta_t=O_P(\sqrt{\eta_t})$ 的含义是：对任意小的 $\epsilon>0$，
可以找到不依赖 $t$ 的阈值 $M$，使归一化误差
$|\delta_t|/\sqrt{\eta_t}$ 超过 $M$ 的概率最终不超过 $\epsilon$。
本题实际上可以得到对所有 $t$ 一致成立的概率界。

因为四次方在非负数上递增，使用 Markov 不等式得到
\begin{align*}
 \Pp(|\delta_t|>M\sqrt{\eta_t})
 &=\Pp(|\delta_t|^4>M^4\eta_t^2)\\
 &\leq\frac{\E|\delta_t|^4}{M^4\eta_t^2}
 \leq\frac{C\eta_t^2}{M^4\eta_t^2}=\frac C{M^4}.
\end{align*}
若 $C>0$，取 $M>(C/\epsilon)^{1/4}$；若 $C=0$，任取 $M>0$。
于是右侧小于 $\epsilon$，所以
\[
 \boxed{\delta_t=O_P(\sqrt{\eta_t}).}
\]

\subsection*{(d) 用可求和概率证明几乎处处收敛}
依概率收敛本身不能直接推出几乎处处收敛。这里利用四阶矩界，
可以证明固定阈值被超过的概率之和有限。
因为 $2\alpha>1$，幂级数判别给出
\[
 \sum_{t=1}^\infty\eta_t^2
 =a^2\sum_{t=1}^\infty(t+t_0)^{-2\alpha}<\infty.
\]
固定正整数 $k$，由 Markov 不等式，
\[
 \Pp(|\delta_t|>1/k)
 \leq k^4\E|\delta_t|^4\leq Ck^4\eta_t^2.
\]
对时间求和可得 $\sum_t\Pp(|\delta_t|>1/k)<\infty$。
第一 Borel--Cantelli 引理说明：这些事件几乎必然只发生有限次。
也就是说，存在概率为 $1$ 的事件 $E_k$，在其上最终有 $|\delta_t|\leq1/k$。
这条引理不要求事件独立。

可数个概率为 $1$ 的事件的交仍然概率为 $1$，所以
$E=\bigcap_{k=1}^\infty E_k$ 满足 $\Pp(E)=1$。
在任意固定的 $\omega\in E$ 上，给定 $\epsilon>0$，选取 $1/k<\epsilon$，
最终便有 $|\delta_t(\omega)|\leq1/k<\epsilon$。这正是逐样本收敛的定义，因此
\[
 \boxed{\delta_t\longrightarrow0\quad\text{几乎处处}.}
\]

\section{凸随机梯度下降的平均误差界}
题设为
\[
 \theta_{t+1}=\theta_t-\eta_tg_t,\qquad
 \E[g_t\mid\mathcal F_t]=\nabla\Phi(\theta_t),\qquad
 \E[\norm{g_t}^2\mid\mathcal F_t]\leq G^2.
\]
$\Phi$ 凸且可微，在 $\theta^\star$ 处达到最小值。记
\[
 e_t=\theta_t-\theta^\star,\qquad
 \Delta_t=\Phi(\theta_t)-\Phi(\theta^\star)\geq0,
 \qquad S_T=\sum_{t=1}^T\eta_t,\quad Q_T=\sum_{t=1}^T\eta_t^2.
\]
所有步长均为确定的正数，故后面的平均权重不是随机量。

\subsection*{预备：确认取期望是合法的}
利用 $\norm{u+v}^2\leq2\norm u^2+2\norm v^2$，
\[
 \E\norm{e_{t+1}}^2
 \leq2\E\norm{e_t}^2+2\eta_t^2\E\norm{g_t}^2
 \leq2\E\norm{e_t}^2+2\eta_t^2G^2.
\]
从 $\E\norm{e_1}^2\leq R^2<\infty$ 出发，归纳得到每个有限时刻的二阶矩有限。
由条件 Jensen 不等式，
\[
 \norm{\nabla\Phi(\theta_t)}^2
 =\norm{\E[g_t\mid\mathcal F_t]}^2
 \leq\E[\norm{g_t}^2\mid\mathcal F_t]\leq G^2.
\]
凸性的一阶不等式给出
$\Delta_t\leq\nabla\Phi(\theta_t)^\top e_t$，再用 Cauchy--Schwarz 得
\[
 0\leq\Delta_t\leq G\norm{e_t},\qquad
 \E\Delta_t\leq G\sqrt{\E\norm{e_t}^2}<\infty.
\]
故函数值误差可积。展开距离时出现的交叉项也可积，因为
$\E|e_t^\top g_t|\leq\sqrt{\E\norm{e_t}^2\E\norm{g_t}^2}<\infty$。

\subsection*{(a) 单步条件期望不等式}
\textbf{第一步：展开距离平方。}
从 $e_{t+1}=e_t-\eta_tg_t$ 出发，
\[
 \norm{e_{t+1}}^2
 =(e_t-\eta_tg_t)^\top(e_t-\eta_tg_t)
 =\norm{e_t}^2-2\eta_te_t^\top g_t+\eta_t^2\norm{g_t}^2.
\]
对 $\mathcal F_t$ 取条件期望，$e_t$ 可移到期望外，得到
\begin{align*}
 \E[\norm{e_{t+1}}^2\mid\mathcal F_t]
 &=\norm{e_t}^2-2\eta_te_t^\top\E[g_t\mid\mathcal F_t]
   +\eta_t^2\E[\norm{g_t}^2\mid\mathcal F_t]\\
 &\leq\norm{e_t}^2-2\eta_te_t^\top\nabla\Phi(\theta_t)+\eta_t^2G^2.
\end{align*}

\textbf{第二步：把梯度内积换成函数值误差。}
凸性在基点 $\theta_t$、目标点 $\theta^\star$ 处给出
\[
 \Phi(\theta^\star)\geq\Phi(\theta_t)
 +\nabla\Phi(\theta_t)^\top(\theta^\star-\theta_t).
\]
移项得 $e_t^\top\nabla\Phi(\theta_t)\geq\Delta_t$。
因为 $-2\eta_t<0$，乘上它以后不等号反向，因此
\[
 \boxed{\E[\norm{e_{t+1}}^2\mid\mathcal F_t]
 \leq\norm{e_t}^2-2\eta_t\Delta_t+\eta_t^2G^2.}
\]
本题只使用普通凸性，不需要强凸性或梯度 Lipschitz 连续性。

\subsection*{(b) 求和并转成加权平均点的误差}
\textbf{第一步：去掉条件期望。}
对 (a) 两边再取期望，并用全期望公式，
\[
 \E\norm{e_{t+1}}^2
 \leq\E\norm{e_t}^2-2\eta_t\E\Delta_t+\eta_t^2G^2.
\]
将误差项移到左侧：
\[
 2\eta_t\E\Delta_t
 \leq\E\norm{e_t}^2-\E\norm{e_{t+1}}^2+\eta_t^2G^2.
\]

\textbf{第二步：对 $t=1,\ldots,T$ 求和。}
相邻距离项逐项抵消，故
\begin{align*}
 2\sum_{t=1}^T\eta_t\E\Delta_t
 &\leq\sum_{t=1}^T\big(\E\norm{e_t}^2-\E\norm{e_{t+1}}^2\big)
   +G^2\sum_{t=1}^T\eta_t^2\\
 &=\E\norm{e_1}^2-\E\norm{e_{T+1}}^2+G^2Q_T\\
 &\leq R^2+G^2Q_T.
\end{align*}
最后一步使用初始二阶矩上界，并去掉一个非正项。

\textbf{第三步：使用凸性处理平均点。}
由于 $S_T>0$，权重 $w_t=\eta_t/S_T$ 满足 $w_t>0$、$\sum_tw_t=1$。
对每条样本路径，凸性保证
\[
 \Phi\left(\overline\theta_T^{(\eta)}\right)
 =\Phi\left(\sum_{t=1}^Tw_t\theta_t\right)
 \leq\sum_{t=1}^Tw_t\Phi(\theta_t).
\]
减去 $\Phi(\theta^\star)=\sum_tw_t\Phi(\theta^\star)$，再取期望，
\begin{align*}
 \E[\Phi(\overline\theta_T^{(\eta)})]-\Phi(\theta^\star)
 &\leq\frac1{S_T}\sum_{t=1}^T\eta_t\E\Delta_t\\
 &\leq\boxed{\frac{R^2+G^2Q_T}{2S_T}}.
\end{align*}
平均点的函数误差非负，并被一个可积的有限加权和控制，
所以这里的期望也存在。

\subsection*{(c) 选取常步长的系数}
题目假设 $R>0$、$G>0$，固定总迭代次数 $T$，取
$\eta_t\equiv\eta=c/\sqrt T$。
这里的步长在这一次长度为 $T$ 的运行中不随 $t$ 改变；它不是 $c/\sqrt t$。

\textbf{第一步：代入步长和与平方步长和。}
\[
 S_T=T\frac c{\sqrt T}=c\sqrt T,
 \qquad Q_T=T\frac{c^2}{T}=c^2.
\]
因为每一步权重相同，
\[
 \overline\theta_T^{(\eta)}
 =\frac1{T\eta}\sum_{t=1}^T\eta\theta_t
 =\frac1T\sum_{t=1}^T\theta_t=\overline\theta_T.
\]
由 (b)，
\[
 \E[\Phi(\overline\theta_T)]-\Phi(\theta^\star)
 \leq\frac{R^2+G^2c^2}{2c\sqrt T}
 =\frac1{2\sqrt T}\left(\frac{R^2}{c}+G^2c\right).
\]

\textbf{第二步：平衡两项，选出最优的保证常数。}
对 $c>0$，直接配方可得
\[
 \frac{R^2}{c}+G^2c-2RG
 =\left(\frac R{\sqrt c}-G\sqrt c\right)^2\geq0.
\]
要使由 (b) 得到的上界达到 $RG/\sqrt T$，令这个平方等于零：
\[
 \frac R{\sqrt c}=G\sqrt c
 \quad\Longleftrightarrow\quad R=Gc
 \quad\Longleftrightarrow\quad\boxed{c=\frac RG}.
\]
$R,G>0$ 保证所选 $c$ 合法。此时 $R^2/c=RG$、$G^2c=RG$，故
\[
 \boxed{\E[\Phi(\overline\theta_T)]-\Phi(\theta^\star)
 \leq\frac{RG+RG}{2\sqrt T}=\frac{RG}{\sqrt T}.}
\]
这个选择使上述统一误差上界最小，并不声称每一个具体问题的实际误差都等于该上界。
\end{document}
